\documentclass[table]{beamer}
%[]中可以使用handout、trancompress等参数

%指定beamer的模式与主题
\mode<presentation>
{
  \usetheme{Madrid}
%\usetheme{Boadilla}
%\usecolortheme{default}
%\usecolortheme{orchid}
%\usecolortheme{whale}
%\usefonttheme{professionalfonts}
}

%\usetheme{Madrid}
%这里还可以选择别的主题：Bergen, Boadilla, Madrid, AnnArbor, CambridgeUS, Pittsburgh, Rochester, Warsaw, ...
%有导航栏的Antibes, JuanLesPins, Montpellier, ...
%有内容的Berkeley, PaloAlto, Goettingen, Marburg, Hannover, ...
%有最小导航栏的Berlin, Ilmenau, Dresden, Darmstadt, Frankfurt, Singapore, Szeged, ...
%有章和节表单的Copenhagen, Luebeck, Malmoe, Warsaw, ...

%\usecolortheme{default}
%设置内部颜色主题(这些主题一般改变block里的颜色);这个主题一般选择动物来命名
%这里还可以选择别的颜色主题,如默认的和有特别目的的颜色主题default,structure,sidebartab,全颜色主题albatross,beetle,crane,dove,fly,seagull,wolverine,beaver

%\usecolortheme{orchid}
%设置外部颜色主题(这些主题一般改变title里的颜色);这个主题一般选择植物来命名
%这里还可以选择别的颜色主题,如默认的和有特别目的的颜色主题lily,orchid,rose

%\usecolortheme{whale}
%设置字体主题;这个主题一般选择海洋动物来命名
%这里还可以选择别的颜色主题,如默认的和有特别目的的颜色主题whale,seahorse,dolphin

%\usefonttheme{professionalfonts}
%类似的还可以定义structurebold,structuresmallcapsserif,professionalfonts


% 控制 beamer 的风格,可以根据自己的爱好修改
%\usepackage{beamerthemesplit} %使用 split 风格
%\usepackage{beamerthemeshadow} %使用 shadow 风格
%\usepackage[width=2cm,dark,tab]{beamerthemesidebar}


% 设定英文字体
\usepackage{fontspec}
%\setmainfont{Times New Roman}
%\setsansfont{Arial}
%\setmonofont{Courier New}

\setsansfont[ItalicFont={Cambria Italic},BoldFont={Cambria Bold},BoldItalicFont={Cambria Bold Italic}]{Cambria}
\setmainfont[ItalicFont={Calibri Italic},BoldFont={Calibri Bold},BoldItalicFont={Calibri Bold Italic}]{Calibri}
\setmonofont[ItalicFont={Consolas Italic},BoldFont={Consolas Bold},BoldItalicFont={Consolas Bold Italic}]{Consolas}


% 设定中文字体
%\usepackage[BoldFont,SlantFont,CJKchecksingle,CJKnumber]{xeCJK}
%\setCJKmainfont[BoldFont={Adobe Heiti Std},ItalicF{SimHei}
%\setCJKmainfont[BoldFont={SimHei},ItalicFont={SimHei}
%\setCJKmainfont{{SimHei}}
%\setCJKsansfont{Adobe Heiti Std}
%\setCJKmonofont{Adobe Fangsong Std}
%\punctstyle{hangmobanjiao}

%\setsansfont{SimHei}
\usepackage[slantfont,boldfont]{xeCJK}

\setCJKmainfont[BoldFont=SimHei,ItalicFont=KaiTi]{SimSun}   % 设置缺省中文字体
\setCJKmonofont[BoldFont=SimHei]{SimSun}   % 设置中文等宽字体
\setCJKsansfont{KaiTi}

\defaultfontfeatures{Mapping=tex-text}
\usepackage{xunicode}
\usepackage{xltxtra}

\XeTeXlinebreaklocale "zh"
\XeTeXlinebreakskip = 0pt plus 1pt minus 0.1pt

\usepackage{setspace}
\usepackage{colortbl,xcolor}
\usepackage{hyperref}
%\hypersetup{xetex,bookmarksnumbered=true,bookmarksopen=true,pdfborder=1,breaklinks,colorlinks,linkcolor=blue,filecolor=black,urlcolor=cyan,citecolor=green}
\hypersetup{xetex,bookmarksnumbered=true,bookmarksopen=true,pdfborder=1,breaklinks,colorlinks,linkcolor=cyan,filecolor=black,urlcolor=blue,citecolor=green}

% 插入图片
\usepackage{graphicx}
% 指定存储图片的路径（当前目录下的figures文件夹）
\graphicspath{{figures/}}

% 可能用到的包
\usepackage{amssymb}
\usepackage{multimedia}
\usepackage{multicol}
\usepackage{subfigure}
\usepackage{graphicx}
\usepackage{amsmath}

% 定义一些自选的模板，包括背景、图标、导航条和页脚等，修改要慎重
% 设置背景渐变由10%的红变成10%的结构颜色
%\beamertemplateshadingbackground{red!10}{structure!10}
%\beamertemplatesolidbackgroundcolor{white!90!blue}
% 使所有隐藏的文本完全透明、动态，而且动态的范围很小
\beamertemplatetransparentcovereddynamic
% 使itemize环境中变成小球，这是一种视觉效果
\beamertemplateballitem
% 为所有已编号的部分设置一个章节目录，并且编号显示成小球
\beamertemplatenumberedballsectiontoc
% 将每一页的要素的要素名设成加粗字体
\beamertemplateboldpartpage

% item逐步显示时,使已经出现的item、正在显示的item、将要出现的item呈现不同颜色
\def\hilite<#1>{
 \temporal<#1>{\color{gray}}{\color{blue}}
    {\color{blue!25}}
}

% 自定义彩色块状结构的颜色
\setbeamercolor{bgcolor}{fg=yellow,bg=cyan}

% 在表格、图片等得标题中显示编号
\setbeamertemplate{caption}[numbered]

% 打开PDF后直接全屏
\hypersetup{pdfpagemode={FullScreen}}



% 使用 \part,\section,\subsection 等命令组织文档结构
% 使用 \frame 命令制作幻灯片

\begin{document}

\logo{\includegraphics[height=0.09\textwidth]{fudan2.png}}
\title[课程项目报告]{基于区域分割的阴影检测和去除算法研究}
\author[刘超颖-14CS]{刘超颖-14CS}
\institute[FDU]{计算机科学与技术学院}
\date{\today}

% 定义目录页
\AtBeginPart{
  \frame{
    \frametitle{\partpage}
    \begin{multicols}{2}
% 如果目录过长，可以打开这个选项分两栏显示
      \tableofcontents
% 使用这个命令自动生成目录
    \end{multicols}
  }  
}

% 在每个Section前都会加入的Frame
\AtBeginSection[]
{
  \begin{frame}<beamer>
    \frametitle{}
	\setcounter{tocdepth}{2}
    \tableofcontents[currentsection,currentsubsection]
  \end{frame}
}
% 在每个Subsection前都会加入的Frame
\AtBeginSubsection[]
{
  \begin{frame}<beamer>
%\begin{frame}<handout:0>
% handout:0 表示只在手稿中出现
    \frametitle{Outline}
	\setcounter{tocdepth}{2}
    \tableofcontents[currentsection,currentsubsection]
% 显示在目录中加亮的当前章节
  \end{frame}
}

\begin{frame}
  \titlepage
\end{frame}

\begin{frame}[plain]
  \frametitle{Table of Contents}
  \setcounter{tocdepth}{2}
  \tableofcontents
\end{frame}


\section{引言}
\begin{frame}
\frametitle{算法的整体思路}
\begin{itemize}
	\setlength\itemsep{1cm}
	\item 传统的全局优化算法难以处理阴影区域内部材质发生变化的情况：
\end{itemize}
\begin{figure}
\center
%\includegraphics[scale=0.2]{1}
\end{figure}

\end{frame}
\begin{frame}
\frametitle{算法的整体思路}
\begin{itemize}
	\item 考虑将图像分割成一块块的小区域；
	\item 以区域为单位进行阴影检测和去除；
	\item 充分利用区域和区域之间是否光照条件相同，是否为相同材料这类信息。
\end{itemize}
\begin{figure}
\begin{minipage}[t]{0.5\linewidth}
\centering
%\includegraphics[width=2.2in]{2}
\end{minipage}%
\begin{minipage}[t]{0.5\linewidth}
\centering
%\includegraphics[width=2.2in]{3}
\end{minipage}
\end{figure}
\end{frame}

\section{阴影检测}
\begin{frame}
\frametitle{算法核心思路}
\begin{itemize}
	\setlength\itemsep{0.5cm}
	\item 对每个区域提特征，判断是否为阴影区；
	\item 区域和区域之间的关系能否为我们判断阴影区域提供帮助？
	\pause
	\begin{itemize}
		\item 相同材料，相同亮度
		\item 相同材料，不同亮度
		\item 不同材料
	\end{itemize}
	\item 构造能量函数，求解最优取值方案.
\begin{align}
\widehat{y}= & \arg \max_y \sum_{i=1}c_i^{shadow}y_i+\alpha_1\sum_{\{i,j\}\in E_{diff}}c_{ij}^{diff}(y_i-y_j) \\
&  -\alpha_2\sum_{\{i,j\}\in E_{same}}c_{ij}^{same}[y_i\ne y_j]
\end{align}
\end{itemize}
\end{frame}
\section{阴影去除}
\begin{frame}
\frametitle{算法核心思路}
\begin{itemize}
	\item 传统方法中阴影区域直接用整个非阴影区域的颜色进行补偿；
	\item 我们将图像分割成区域后，可以单独考虑每块阴影区域的颜色；
	\item 如果一块阴影区与一块非阴影区具有相同的材质，则按照这块非阴影区的颜色去修改阴影区比较合适；
	\item CIEL*a*b空间直方图均衡；
\end{itemize}
\begin{figure}
\center
%\includegraphics[scale=0.3]{10}
\end{figure}
\end{frame}
\begin{frame}
\frametitle{算法的改进}
\begin{itemize}
	\item 原算法对边缘的处理不够充分：
	\begin{itemize}
		\item 存在一些半影区的像素没有被标记为阴影像素；
		\item 在非阴影区内部可能存在小于一半的阴影像素；
	\end{itemize}
\end{itemize}
\begin{figure}
\center
%\includegraphics[scale=0.3]{4}
\end{figure}
\end{frame}
\begin{frame}
\frametitle{算法的改进}
\begin{itemize}
	\item 使用形态学方法找出这些可能是阴影或者未处理的阴影像素；
	\item 用in-painting算法对这些区域进行填充；
\end{itemize}
\begin{figure}
\center
%\includegraphics[scale=0.3]{5}
\end{figure}
\end{frame}
\section{实验}
\begin{frame}
\frametitle{算法的实现}
\begin{itemize}
	\item 阴影检测和阴影去除的核心代码均用Matlab实现；
	\item 比较费时的过程用C++实现再封装为Matlab借口，比如SAP最大流算法；
	\item 为了进行对比实现了一种传统的全局处理的方法；
\end{itemize}
\end{frame}
\begin{frame}
\frametitle{阴影检测算法}
\begin{itemize}
	\item 使用了三个数据集进行测试，评价方法采用逐像素比较；
\end{itemize}
\begin{figure}
\center
%\includegraphics[scale=0.4]{6}
\end{figure}
\end{frame}
\begin{frame}
\frametitle{阴影去除算法}
\begin{itemize}
	\item 使用RMSE评价方法：
\end{itemize}
\begin{figure}
\center
%\includegraphics[scale=0.5]{9}
\end{figure}
\end{frame}
\begin{frame}
\frametitle{实际运行效果}
%\includegraphics[scale=0.45]{7}
\end{frame}
\begin{frame}
\frametitle{实际运行效果}
\begin{figure}
\begin{minipage}[t]{0.5\linewidth}
\centering
%\includegraphics[width=2.2in]{11}
\end{minipage}%
\begin{minipage}[t]{0.5\linewidth}
\centering
%\includegraphics[width=2.2in]{12}
\end{minipage}
\end{figure}
\end{frame}
\begin{frame}
	\begin{spacing}{1.5}
	\begin{center}
	\Huge{\textbf{Thanks for your attention!}}

	\Huge{\textit{Any questions?}}
\end{center}
\end{spacing}
\end{frame}
\end{document}
